\scanpage{145}
\margin{L'indice $d_1$ plus haut devient ici $d_0$.}
\inserted{de codim $d$,} qui se factorise par $\Theta^{(d)}$ (NB il est essentiel pour ce qu'on veut faire qu'on n'exige \uline{pas} que $D'\subset\Theta^{(d)*}$), \deleted{\unsure{on en}} d'où un ``voisinage tubulaire'' $\mathcal V_{D',Y}^{\unsure{\deleted{*}}}$ de $D'$ dans $Y$, et par trace de $Y^{*}=Y\setminus\Theta$ un \deleted{\unsure{un}} germe d'espace analytique $\mathcal V^{*}_{D',Y}$. Alors on a une équivalence de groupoïdes, \uline{\unsure{définie} à \deleted{équivalence} isomorphisme unique près}
\[
(239)\qquad\Pi_1\mathcal V^{*}_{D',Y}\xleftarrow{\sim}\Pi_1P_{D'^{*}}
\]
\uline{Remarque} Le cas le plus simple est celui où $Y$ est de dim $1$, et $d=1$, $D'$ réduit à un pt --- le \deleted{\unsure{cas}} cas peut paraître trivial, néanmoins la question de canonicité de l'équivalence m'a semblé nécessiter un soin certain, et fait l'objet de la longue digression des ``notes \unsure{Marche}'' § 53 n° 50 (p. \unsure{768'|7681} ff). Ici \inserted{(en général)} il devrait être ``assez évident'' de définir \inserted{(de façon non canonique)} un plongement \deleted{topologique} \inserted{\unsure{cl.} $C^{\infty}$} (non analytique !)
\margin{144}
